% !TeX program = pdflatex
% ======================================================================
% DOCUMENTO MAESTRO: ARTICULO REGULAR PARA IEEE T-IV O IEEE TVT
%
% Guardar como master.tex en UTF-8 y compilar dos veces con pdfLaTeX.
% Requiere IEEEtran y paquetes habituales de TeX Live/MiKTeX/Overleaf.
% Bibliografía integrada: no requiere archivos .bib ni imágenes externas.
%
% BASE EDITORIAL Y TRAZABILIDAD
% Consulta de normas y biblioteca: 15 de octubre de 2026.
%
% Fuentes oficiales:
% T-IV:
% https://ieee-itss.org/pub/t-iv/
% TVT:
% https://vtsociety.org/publication/ieee-transactions-vehicular-technology/guidelines-authors/instructions
% Estructura general IEEE:
% https://journals.ieeeauthorcenter.ieee.org/create-your-ieee-journal-article/create-the-text-of-your-article/structure-your-article/
% Selector oficial de plantillas:
% https://template-selector.ieee.org/
%
% Biblioteca Zotero consultada en modo lectura:
% Colección «Articulo IEEET», clave TAW68QTE.
%
% Ejemplos recientes utilizados para orientar esta plantilla:
% 1. Bae et al., T-IV, 2025:
%    costes programados, aprendizaje de expertos y evaluación comparativa.
%    Zotero: TB7VGR26. DOI: 10.1109/TIV.2024.3465598.
% 2. Zheng et al., T-IV, 2026:
%    estabilidad, adaptación del controlador y validación HIL.
%    Zotero: JQZK5CRG. DOI: 10.1109/TIV.2024.3398129.
% 3. Chen et al., TVT, 2026:
%    modelos informados por física y validación en vehículo real.
%    Zotero: ZZ2QE4FS. DOI: 10.1109/TVT.2025.3635878.
% 4. Yang et al., TVT, 2025:
%    control basado en datos, factibilidad recursiva y estabilidad.
%    Zotero: VEDMJ92F. DOI: 10.1109/TVT.2025.3552933.
% 5. Li et al., TVT, 2025:
%    modelado híbrido y adelantamientos con diferentes condiciones.
%    Zotero: XFFMBW52. DOI: 10.1109/TVT.2024.3453170.
%
% Se conserva el año del volumen registrado en Zotero; puede diferir
% del año de publicación anticipada y del año incluido en el DOI.
%
% ALCANCE DE LA ESTRUCTURA
% Las revistas no imponen una jerarquía universal de títulos y subtítulos.
% Este índice es una propuesta editorial para control vehicular con
% aprendizaje y MPC, basada en las normas IEEE y los resúmenes consultados.
% No se verificaron los índices íntegros de los artículos: los encabezados
% siguientes no se presentan como una transcripción literal de ellos.
% Adaptar los títulos al problema y a las contribuciones reales.
%
% Los porcentajes y números de párrafos son recomendaciones editoriales,
% no límites oficiales ni criterios particulares de revisores identificados.
%
% REQUISITOS DE PRESENTACION
% T-IV: 10 páginas sugeridas para artículo regular; admite sobrelongitud.
% TVT: máximo inicial de 14 páginas; revisado/final de 16 páginas,
% incluidas referencias y biografías. Dos columnas y fuente mínima de 10 pt.
% Ambas revistas mantienen visible la identidad de los autores.
% Texto publicable en inglés; comentarios de orientación en español.
% Mantener los márgenes y espaciados establecidos por IEEEtran.
% ======================================================================

\documentclass[10pt,journal]{IEEEtran}

% -------------------------- CONFIGURACION -----------------------------

\newif\ifTargetTIV
\TargetTIVtrue
% Cambiar a \TargetTIVfalse para IEEE Transactions on Vehicular Technology.

\newif\ifIncludeTheory
\IncludeTheorytrue
% Desactivar si el artículo no contiene una contribución teórica propia.

\newif\ifIncludeAppendix
\IncludeAppendixfalse
% Activar para demostraciones o derivaciones extensas.

\newif\ifIncludeDeclarations
\IncludeDeclarationsfalse
% Activar cuando proceda incluir declaraciones dentro del manuscrito.
% Completar también los formularios editoriales correspondientes.

\newif\ifIncludeBiographies
\IncludeBiographiestrue
% Mantener las biografías para el artículo regular de T-IV.

\ifTargetTIV
  \newcommand{\TargetJournal}{IEEE Transactions on Intelligent Vehicles}
\else
  \newcommand{\TargetJournal}{IEEE Transactions on Vehicular Technology}
\fi

% ------------------------------ PAQUETES ------------------------------

\usepackage[T1]{fontenc}
\usepackage[utf8]{inputenc}
\usepackage[english]{babel}

\usepackage{amsmath,amssymb,amsfonts}
\usepackage{bm}
\interdisplaylinepenalty=2500

\usepackage{graphicx}
\usepackage{booktabs,array,tabularx}
\usepackage{siunitx}
\sisetup{detect-all}

\usepackage{algorithmic}
% El pseudocódigo puede presentarse dentro de figure, evitando introducir
% un flotante algorithm con convenciones distintas de las de IEEE.

\usepackage[numbers,sort&compress]{natbib}
% Citas numéricas. No cargar simultáneamente cite ni biblatex.
% Si se utiliza BibTeX externo con natbib, emplear IEEEtranN.

\usepackage{url}
\usepackage[hidelinks]{hyperref}

\hypersetup{
  pdftitle={Master Template for an IEEE Vehicular Control Article},
  pdfsubject={\TargetJournal},
  pdfauthor={Replace with the manuscript authors}
}

% No modificar el formato mediante geometry, titlesec o setspace.
% Para subfiguras, si fueran necesarias:
% \usepackage[caption=false,font=footnotesize]{subfig}

% ---------------------------- HERRAMIENTAS ----------------------------

\newcommand{\placeholder}[1]{\textit{[#1]}}
\newcommand{\R}{\mathbb{R}}
\newcommand{\vect}[1]{\bm{#1}}
\newcommand{\norm}[1]{\left\lVert#1\right\rVert}
\newcommand{\argmin}{\operatorname*{arg\,min}}

\newtheorem{assumption}{Assumption}
\newtheorem{proposition}{Proposition}
\newtheorem{theorem}{Theorem}
\newtheorem{lemma}{Lemma}

\begin{document}

% TITULO
% Contenido: método, tarea vehicular y condición distintiva.
% Tono/extensión: específico y conciso; unas 12--18 palabras orientativas.
% Revisión: evitar reclamos de optimalidad o seguridad no demostrados.

\title{[Method] for [Vehicle Control Task] Under [Operating Conditions]}

% AUTORES
% Nombres bajo el título; afiliaciones y correos en la primera nota.
% Sustituir todos los marcadores. Añadir membresías IEEE solo si son ciertas.

\author{First A. Author, Second B. Author, and Third C. Author%
\thanks{First A. Author and Second B. Author are with [Department],
[Institution], [City], [Country] (e-mail: [institutional addresses]).}%
\thanks{Third C. Author is with [Department], [Institution], [City],
[Country] (e-mail: [institutional address]).}%
\thanks{Corresponding author: [Full name].
This work was supported by [Funder] under Grant [Identifier].}}

% Verificar la financiación; retirar o adaptar la oración si corresponde.
% No inventar fechas de recepción/aceptación, DOI, volumen o copyright.

\maketitle

\begin{abstract}
% CONTENIDO: problema concreto, brecha, mecanismo propuesto, evaluación,
% resultados cuantitativos reales y alcance de la conclusión.
% TONO/EXTENSION: un párrafo de 150--250 palabras. Distribución orientativa:
% 2 frases de problema, 2--3 de método y 2 de evidencia y alcance.
% REVISION: novedad reconocible, cifras contextualizadas y correspondencia
% con los resultados. Evitar citas, abreviaturas y ecuaciones desplegadas.
% Redactar al terminar el artículo. El marcador no es un resumen final.
\placeholder{Write the problem, research gap, proposed mechanism,
evaluation protocol, measured results, and supported conclusion here.}
\end{abstract}

\begin{IEEEkeywords}
% CONTENIDO: términos que representen sistema, tarea y método.
% TONO/EXTENSION: 3--5 términos específicos, preferiblemente alfabéticos.
% REVISION: utilizar el tesauro IEEE cuando corresponda y evitar etiquetas
% que no estén sustentadas por el contenido del artículo.
Automated vehicles, learning-based control, model predictive control,
trajectory tracking.
\end{IEEEkeywords}

% ======================================================================
\section{Introduction}
\label{sec:introduction}
% CONTENIDO: motivación vehicular, literatura crítica, brecha y aportes.
% TONO/EXTENSION: 6--9 párrafos y lista breve; aproximadamente 15% del cuerpo.
% REVISION: pertinencia para la revista, diferencia frente al trabajo más
% próximo y correspondencia entre cada contribución y su evidencia.

\subsection{Background and Motivation}
% CONTENIDO: tarea, dominio operativo y consecuencias de la limitación actual.
% TONO/EXTENSION: 1--2 párrafos, de lo general a lo específico.
% REVISION: problema medible. Para T-IV, destacar inteligencia y operación
% automatizada; para TVT, precisar el avance tecnológico vehicular.
\placeholder{Describe the vehicle task, operating conditions, practical
consequences, and unresolved technical challenge.}

\subsection{Related Work and Research Gap}
% CONTENIDO: organizar la literatura por mecanismos y limitaciones.
% TONO/EXTENSION: 3--4 párrafos de síntesis crítica; evitar una lista de autores.
% REVISION: trabajos recientes y seminales, comparación con alternativas
% cercanas y una brecha verificable que conduzca al objetivo del artículo.
% Si amplía una publicación de congreso propia, citarla y explicar el avance.
%
% Las citas siguientes son anclas de Zotero. Conservar solo las pertinentes
% después de leer los artículos completos y verificar las afirmaciones.

\placeholder{Compare learning-based predictive models and hierarchical
vehicle control approaches.}%
\cite{liDataDrivenHierarchicalModel2025,chenExperimentalEvaluationTrajectory2026a}

\placeholder{Contrast adaptive objectives, stability-aware control, and
data-driven feasibility analyses; identify the exact remaining gap.}%
\cite{baeLearningExpertLevelRacing2025,zhengAdaptiveModelPredictive2026,
yangDataDrivenModelPredictive2025}

\subsection{Main Contributions}
% CONTENIDO: 2--3 aportes distintos: avance, diferencia y evidencia.
% TONO/EXTENSION: una frase introductoria y viñetas de 2--3 líneas.
% REVISION: separar novedad metodológica, propiedad analítica y validación.
% La combinación de técnicas conocidas requiere una ventaja demostrable.
\placeholder{State the research objective and list the verified contributions.}

\begin{itemize}
  \item \placeholder{Methodological advance and difference from the
  closest prior method; identify the supporting section.}
  \item \placeholder{Analytical or computational advance and its
  assumptions; retain only if established in this study.}
  \item \placeholder{Experimental evidence that resolves the research
  question, including operating conditions and reproducibility.}
\end{itemize}

% ORGANIZACION: opcional; un párrafo de 2--3 frases, sin subtítulo propio.
% Mantener las referencias automáticas al modificar la estructura.
The remainder of this article presents the problem formulation in
Section~\ref{sec:problem} and the proposed method in
Section~\ref{sec:method}.
\ifIncludeTheory
Section~\ref{sec:theory} develops the theoretical analysis.
\fi
Sections~\ref{sec:setup} and~\ref{sec:results} describe the evaluation
and discuss the results. Section~\ref{sec:conclusion} concludes the article.

% ======================================================================
\section{System Modeling and Problem Formulation}
\label{sec:problem}
% CONTENIDO: sistema, variables, hipótesis, incertidumbre y objetivo formal.
% TONO/EXTENSION: aproximadamente 12% del cuerpo; exposición reproducible.
% REVISION: dimensiones, unidades, información disponible y coherencia
% entre simplificaciones del modelo y condiciones experimentales.

\subsection{Vehicle Model and Notation}
% CONTENIDO: estados, entradas, salidas, dinámica, muestreo y coordenadas.
% TONO/EXTENSION: 2--3 párrafos y las ecuaciones indispensables.
% REVISION: distinguir planta, predictor y estimador; identificar parámetros
% medidos o identificados y definir cada variable antes de utilizarla.
\placeholder{Define the vehicle dynamics, measurements, control inputs,
coordinate frames, sampling period, and model parameters.}

% ECUACION DE PRUEBA: descomentar y adaptar.
% Definir x_k en R^n, u_k en R^m y w_k en R^n, con unidades por componente.
% Explicar la discretización y el significado de las perturbaciones.
%
% \begin{equation}
%   \vect{x}_{k+1}
%   = f_{\theta}(\vect{x}_k,\vect{u}_k)+\vect{w}_k .
%   \label{eq:dynamics}
% \end{equation}
%
% La ecuación forma parte de una oración y necesita puntuación.
% Referenciar mediante \eqref{eq:dynamics} después de activarla.

\subsection{Operating Domain, Uncertainty, and Assumptions}
% CONTENIDO: velocidades, fricción, carga, sensores, actuadores y entorno.
% TONO/EXTENSION: 1--2 párrafos; enumerar solo supuestos utilizados después.
% REVISION: ruido, retardos, saturación, variables observables y dominio
% de validez. Justificar las cotas utilizadas con datos o referencias.
\placeholder{Specify the operating domain, measurable quantities,
disturbance model, actuator limits, and justified assumptions.}

% \begin{assumption}\label{ass:disturbance}
%   [State a precise, testable assumption and its domain of validity.]
% \end{assumption}

\subsection{Control Objectives and Constraints}
% CONTENIDO: tarea, referencia, restricciones duras/blandas y criterios de fallo.
% TONO/EXTENSION: 1--2 párrafos y formulación matemática compacta.
% REVISION: objetivos en tensión, factibilidad y disponibilidad de referencia;
% distinguir satisfacción de restricciones y garantías de seguridad.
\placeholder{Formulate the control problem, tracking objectives,
safety-related constraints, and admissible operating conditions.}

% ======================================================================
\section{Proposed Learning-Based Predictive Control Method}
\label{sec:method}
% CONTENIDO: explicar el mecanismo desde las mediciones hasta la acción.
% TONO/EXTENSION: aproximadamente 25% del cuerpo; técnico y reproducible.
% REVISION: novedad localizada, integración aprendizaje--control y detalle
% suficiente para implementar el método sin consultar código privado.
% Adaptar el título si el aporte real no utiliza aprendizaje o MPC.

\subsection{Overall Architecture}
% CONTENIDO: bloques, interfaces, flujo de información y escalas temporales.
% TONO/EXTENSION: 1--2 párrafos y un diagrama de arquitectura.
% REVISION: diferenciar entrenamiento, inferencia, optimización y planta;
% señalar dónde se introduce la contribución.
\placeholder{Describe the architecture, module interfaces, timing,
and the location of the proposed technical contribution.}

% FIGURA DE PRUEBA: descomentar el entorno completo.
% El recuadro permite compilar sin un archivo gráfico externo.
% Para una figura real, sustituir el recuadro por:
% \includegraphics[width=\columnwidth]{figures/architecture.pdf}
%
% \begin{figure}[!t]
%   \centering
%   \fbox{\parbox[c][24mm][c]{0.90\columnwidth}{\centering
%     Measurements $\rightarrow$ predictor $\rightarrow$ controller\\
%     $\uparrow$\hspace{35mm}$\downarrow$\\
%     Vehicle and environment}}
%   \caption{Architecture of the proposed control system. Solid arrows
%   denote online information flow; identify offline training separately.}
%   \label{fig:architecture}
% \end{figure}
%
% Pie debajo; label después de caption. Usar figure* para dos columnas.
% Preferir gráficos vectoriales y tipografía legible al tamaño definitivo.
% Citar la figura en el texto antes de analizar sus detalles.

\subsection{Learning-Based Prediction Model}
% CONTENIDO: qué se aprende y cómo se integra con el modelo nominal.
% TONO/EXTENSION: 2--4 párrafos distribuidos entre los subapartados.
% REVISION: causalidad temporal, cobertura de datos y tratamiento del error.
\placeholder{Define the learned quantity and its interface with the
nominal dynamics and the predictive controller.}

\subsubsection{Offline Identification and Training}
% CONTENIDO: entradas/salidas, arquitectura, pérdida, regularización,
% normalización, optimizador y selección de hiperparámetros.
% TONO/EXTENSION: 1--2 párrafos y una ecuación si aclara el aporte.
% REVISION: particiones por trayectorias, vehículos o sesiones; prevenir
% fuga temporal y selección de hiperparámetros con datos de prueba.
\placeholder{Specify identification data, training procedure,
hyperparameters, and the validation-only selection protocol.}

\subsubsection{Online Inference and Uncertainty Handling}
% CONTENIDO: ventana histórica, incertidumbre, actualización y latencia.
% TONO/EXTENSION: 1--2 párrafos; precisar si el modelo permanece congelado.
% REVISION: información disponible durante la ejecución y comportamiento
% fuera de distribución; distinguir inferencia de aprendizaje en línea.
\placeholder{Describe online inputs, uncertainty treatment,
and any adaptation mechanism or frozen-model assumption.}

\subsection{Predictive Control Formulation}
% CONTENIDO: optimización completa, horizonte y política de aplicación.
% TONO/EXTENSION: 2--3 párrafos, ecuaciones alineadas y símbolos definidos.
% REVISION: coherencia entre predictor, coste, restricciones y solver.
\placeholder{Present the complete finite-horizon optimization problem
and explain the receding-horizon implementation.}

\subsubsection{Objective Function}
% CONTENIDO: seguimiento, esfuerzo, término terminal y objetivos adicionales
% pertinentes, como confort o energía; definir pesos y escalas.
% TONO/EXTENSION: un párrafo y una ecuación compacta.
% REVISION: justificar la sintonía sin optimizar sobre el conjunto de prueba.
\placeholder{Define the stage and terminal costs and justify their weights.}

\subsubsection{Constraints and Feasibility Management}
% CONTENIDO: límites físicos, restricciones vehiculares, holguras y terminal.
% TONO/EXTENSION: 1--2 párrafos y restricciones explícitas.
% REVISION: viabilidad inicial, penalización de holguras e infactibilidad;
% evitar presentar restricciones blandas como garantías duras.
\placeholder{Define admissible states and inputs, slack penalties,
terminal ingredients, and the infeasibility handling policy.}

\subsection{Implementation and Computational Complexity}
% CONTENIDO: solver, tolerancias, inicialización, parada y acción de respaldo.
% TONO/EXTENSION: 1--2 párrafos y pseudocódigo si elimina ambigüedades.
% REVISION: tamaño del problema y coste por paso; diferenciar complejidad
% asintótica y tiempos medidos en una plataforma concreta.
\placeholder{Provide a reproducible execution procedure, solver settings,
computational scaling, and timeout handling.}

% PSEUDOCODIGO OPCIONAL:
%
% \begin{figure}[!t]
%   \begin{algorithmic}[1]
%     \REQUIRE Measurement, reference, model, and previous solution
%     \STATE Update the state estimate and predictive model
%     \STATE Construct and solve the constrained control problem
%     \IF{the solution satisfies feasibility and timing requirements}
%       \STATE Apply the first optimized control input
%     \ELSE
%       \STATE Apply the validated fallback action
%     \ENDIF
%   \end{algorithmic}
%   \caption{Online control procedure.}
%   \label{fig:procedure}
% \end{figure}

% ======================================================================
\ifIncludeTheory
\section{Theoretical Analysis}
\label{sec:theory}
% CONTENIDO: propiedades demostrables del método propuesto.
% TONO/EXTENSION: aproximadamente 8% del cuerpo; precisión matemática.
% REVISION: supuestos explícitos y conexión con la implementación real.
% Sección condicional: no es una exigencia universal de T-IV o TVT.
% Si se elimina, reasignar el espacio al método o a la evaluación.

\subsection{Recursive Feasibility and Constraint Satisfaction}
% CONTENIDO: condiciones iniciales, conjunto terminal y tratamiento del error.
% TONO/EXTENSION: 1--2 párrafos, proposición y demostración.
% REVISION: construcción de solución candidata y alcance determinista
% o probabilístico. Ensayos sin fallos no prueban factibilidad recursiva.
\placeholder{State and prove the applicable feasibility property,
including all assumptions and its relationship to implementation.}

\subsection{Closed-Loop Stability or Error Bounds}
% CONTENIDO: estabilidad, cota de error o probabilidad de violación.
% TONO/EXTENSION: 1--2 párrafos, resultado y demostración o remisión al apéndice.
% REVISION: efecto del error de aproximación y perturbaciones; distinguir
% convergencia del entrenamiento y estabilidad del sistema controlado.
% Renombrar este subtítulo según la propiedad realmente establecida.
\placeholder{Establish the claimed closed-loop property and explain
the assumptions under which it holds.}

% \begin{theorem}\label{thm:closed-loop}
%   [Insert a precise statement supported by the actual analysis.]
% \end{theorem}
%
% \begin{IEEEproof}
%   [Give the proof or refer to a complete derivation in the appendix.]
% \end{IEEEproof}
\fi

% ======================================================================
\section{Experimental Setup}
\label{sec:setup}
% CONTENIDO: protocolo definido antes de observar los resultados principales.
% TONO/EXTENSION: aproximadamente 10% del cuerpo; factual y reproducible.
% REVISION: comparación justa, cobertura del dominio y detalle de plataforma.

\subsection{Platforms and Operating Scenarios}
% CONTENIDO: simulador/vehículo, sensores, hardware, software, muestreo,
% trayectorias, velocidades, fricción, carga y perturbaciones evaluadas.
% TONO/EXTENSION: 1--2 párrafos y tabla de parámetros cuando sea útil.
% REVISION: distinguir simulación, HIL y vehículo real; definir criterios
% de fallo y no atribuir a una modalidad evidencia obtenida en otra.
\placeholder{Describe the platforms, scenario matrix, parameter ranges,
execution hardware, and experimental termination criteria.}

\subsection{Data Protocol and Baseline Methods}
% CONTENIDO: procedencia y particiones de datos, ensayos y comparadores;
% incluir una referencia convencional fuerte y alternativas recientes.
% TONO/EXTENSION: 2 párrafos y configuración explícita o tabla compacta.
% REVISION: mismos datos, restricciones e información; presupuesto de
% sintonía equitativo y documentación de código, versiones y modificaciones.
\placeholder{Specify independent data partitions, trial counts, random
seeds, baseline implementations, and fair tuning conditions.}

\subsection{Evaluation Metrics and Statistical Protocol}
% CONTENIDO: métricas de tarea, restricciones, confort/energía si proceden,
% tiempos, fallos, unidad de réplica e incertidumbre de estimación.
% TONO/EXTENSION: 1--2 párrafos; definir las métricas no estándar.
% REVISION: tamaños de efecto y dispersión. No tratar muestras temporales
% correladas como ensayos independientes; justificar tamaño muestral y
% pruebas estadísticas cuando se utilicen.
\placeholder{Define primary and secondary metrics, aggregation rules,
independent repetitions, confidence intervals, and failure reporting.}

% TABLA PROFESIONAL DE PRUEBA:
% Descomentar el entorno completo y sustituir -- por resultados medidos.
% Definir si se reporta media/desviación o mediana/rango intercuartil.
%
% \begin{table}[!t]
%   \caption{Closed-loop performance under the same evaluation protocol}
%   \label{tab:comparison}
%   \centering
%   \footnotesize
%   \begin{tabular}{@{}lccc@{}}
%     \toprule
%     Method & RMSE (\si{\metre}) & Time (\si{\milli\second})
%            & Failures (\%) \\
%     \midrule
%     Conventional MPC & -- & -- & -- \\
%     Recent baseline  & -- & -- & -- \\
%     Proposed method  & -- & -- & -- \\
%     \bottomrule
%   \end{tabular}
%   \par\smallskip
%   \parbox{\columnwidth}{\footnotesize
%     Report the number of independent runs and the uncertainty summary.
%     Define failure and whether smaller values are preferable.}
% \end{table}
%
% Título encima; evitar líneas verticales y reducción ilegible de fuente.
% No sugerir superioridad estadística con negritas sin análisis suficiente.

% ======================================================================
\section{Results and Discussion}
\label{sec:results}
% CONTENIDO: responder las preguntas de investigación y explicar mecanismos.
% TONO/EXTENSION: aproximadamente 25% del cuerpo; evidencia e interpretación.
% REVISION: comparaciones sólidas, casos adversos y alcance de las conclusiones.
% Separar Discussion como sección propia si la extensión lo justifica.

\subsection{Prediction Accuracy and Model Validation}
% CONTENIDO: error fuera de muestra, horizontes y regímenes relevantes.
% TONO/EXTENSION: 1--2 párrafos y una figura/tabla si hay aporte de modelado.
% REVISION: comparar con el modelo nominal y un predictor alternativo fuerte;
% comprobar si una mejor predicción produce una mejora de control.
\placeholder{Report held-out prediction performance and relate modeling
errors to the downstream control task.}

\subsection{Closed-Loop Performance and Constraint Satisfaction}
% CONTENIDO: métricas principales, restricciones, fallos y condiciones difíciles.
% TONO/EXTENSION: 2--3 párrafos, tabla agregada y trayectorias representativas.
% REVISION: magnitud y variabilidad de las mejoras, costes asociados y
% resultados desfavorables; evitar seleccionar únicamente casos favorables.
\placeholder{Compare the methods across the full scenario matrix,
including performance, constraint violations, and unsuccessful trials.}

\subsection{Ablation and Sensitivity Analysis}
% CONTENIDO: retirar mecanismos nuevos y variar hiperparámetros críticos.
% TONO/EXTENSION: 1--2 párrafos y tabla/gráfica centrada en aislar efectos.
% REVISION: contribución individual de cada componente y control de recursos;
% no confundir un mejor algoritmo con un mayor presupuesto computacional.
\placeholder{Isolate the contributions of the proposed components and
report sensitivity to the parameters that affect deployment.}

\subsection{Real-Time Performance and Experimental Validation}
% CONTENIDO: latencia completa, hardware, distribución de tiempos,
% retrasos, timeouts y respaldo; pruebas HIL o reales si están disponibles.
% TONO/EXTENSION: 1--2 párrafos con condiciones y medidas completas.
% REVISION: media, percentiles altos y máximo observado frente al muestreo;
% diferenciar un máximo observado y una cota demostrada de peor caso.
\placeholder{Report implementation timing, missed deadlines, and the
available hardware-in-the-loop or real-vehicle validation.}

\subsection{Limitations and Generalization}
% CONTENIDO: fallos, supuestos restrictivos y transferencia a otros vehículos,
% carreteras, sensores o dominios operativos.
% TONO/EXTENSION: 1--2 párrafos analíticos y específicos.
% REVISION: conclusiones proporcionales a la evidencia; distinguir resultados
% comprobados, hipótesis de transferencia y cuestiones todavía abiertas.
\placeholder{Discuss the observed limitations, unresolved failure modes,
generalization evidence, and practical deployment implications.}

% ======================================================================
\section{Conclusion}
\label{sec:conclusion}
% CONTENIDO: problema abordado, mecanismo distintivo, hallazgo principal
% cuantificado y siguiente paso derivado de una limitación concreta.
% TONO/EXTENSION: 1--2 párrafos; aproximadamente 5% del cuerpo.
% REVISION: responder a los aportes iniciales sin incorporar experimentos,
% cifras o afirmaciones nuevas ni ampliar injustificadamente el alcance.
\placeholder{Summarize the verified contribution, the most informative
measured findings, the supported scope, and a concrete next research step.}

% ======================================================================
\ifIncludeAppendix
\appendices

\section{Additional Derivations}
\label{app:proofs}
% CONTENIDO: pasos técnicos extensos necesarios para verificar los resultados.
% TONO/EXTENSION: lo indispensable; cuenta en la extensión del manuscrito.
% REVISION: pruebas completas y trazables; mantener en el cuerpo los supuestos
% y resultados principales.
\placeholder{Provide the complete supporting derivations referenced
from the main text.}
\fi

% ======================================================================
\section*{Acknowledgment}
% CONTENIDO: asistencia y apoyos verificables; financiación en la primera
% nota de autor según el formato IEEE, evitando duplicaciones innecesarias.
% TONO/EXTENSION: 2--4 frases; nombres y reconocimientos confirmados.
% REVISION: exactitud de los apoyos y declaraciones editoriales aplicables.
% TVT solicita reconocer el uso de IA para modificar texto de los autores
% cuando se haya utilizado; describir herramienta y alcance reales.
% Esta plantilla contiene estructura y marcadores, no resultados científicos.
\placeholder{Insert verified acknowledgments and any applicable disclosure.}

\ifIncludeDeclarations
\section*{Conflict of Interest}
% CONTENIDO: declaración verificada por todos los autores.
% TONO/EXTENSION: 1--2 frases; no afirmar ausencia de conflictos por defecto.
% REVISION: consistencia con los formularios editoriales. T-IV solicita
% declarar COI en su checklist; no implica exigir este encabezado impreso.
\placeholder{Insert the authors' verified competing-interest statement.}

\section*{Data and Code Availability}
% CONTENIDO: repositorio real, versión, licencia y restricciones de acceso.
% TONO/EXTENSION: 1--3 frases, con enlace persistente cuando exista.
% REVISION: correspondencia entre artefactos y resultados, y posibilidad de
% reproducirlos. No es una sección obligatoria universal de ambas revistas.
\placeholder{State the actual availability of data, code, configurations,
and instructions, including any access restrictions.}
\fi

% ======================================================================
% REFERENCIAS
% Orden numérico por primera aparición y formato IEEE.
% Entradas procedentes de Zotero; revisar pertinencia antes del envío.
% Las claves de cite corresponden a citationKey, no al identificador
% de ocho caracteres del ítem Zotero.
%
% Para migrar a una bibliografía externa exportada desde Zotero:
% 1. Eliminar el entorno thebibliography completo.
% 2. Incorporar:
%    \bibliographystyle{IEEEtranN}
%    \bibliography{references}
% 3. Compilar: pdfLaTeX, BibTeX, pdfLaTeX, pdfLaTeX.

\begin{thebibliography}{00}

\bibitem{liDataDrivenHierarchicalModel2025}
Y.~Li, Z.~Chen, T.~Wang, X.~Zeng, and Z.~Yin,
``Data-driven hierarchical model predictive control for automated
overtaking maneuver via Gaussian process regression,''
\emph{IEEE Trans. Veh. Technol.}, vol.~74, no.~1,
pp.~263--278, Jan. 2025, doi:
\href{https://doi.org/10.1109/TVT.2024.3453170}{10.1109/TVT.2024.3453170}.

\bibitem{chenExperimentalEvaluationTrajectory2026a}
G.~Chen, X.~Zhao, Z.~Gao, Z.~Gao, X.~Wang, and M.~Hua,
``Experimental evaluation of trajectory tracking with nonlinear model
predictive control enhanced by physics-informed neural network,''
\emph{IEEE Trans. Veh. Technol.}, vol.~75, no.~6,
pp.~8928--8939, Jun. 2026, doi:
\href{https://doi.org/10.1109/TVT.2025.3635878}{10.1109/TVT.2025.3635878}.

\bibitem{baeLearningExpertLevelRacing2025}
J.~Bae, J.~Lim, B.~Suh, J.~Lee, K.~Ryu, J.~Kim, and J.~Choi,
``Learning expert-level racing strategies via scheduled cost functions
in model predictive control,''
\emph{IEEE Trans. Intell. Veh.}, vol.~10, no.~6,
pp.~3827--3840, Jun. 2025, doi:
\href{https://doi.org/10.1109/TIV.2024.3465598}{10.1109/TIV.2024.3465598}.

\bibitem{zhengAdaptiveModelPredictive2026}
Z.~Zheng, S.~Wang, X.~Zhao, Q.~Yu, H.~Zhang, Y.~Lv, and J.~Tian,
``Adaptive model predictive control of four-wheel drive intelligent
electric vehicles based on stability probability under extreme
braking conditions,''
\emph{IEEE Trans. Intell. Veh.}, vol.~11, no.~5,
pp.~618--628, May 2026, doi:
\href{https://doi.org/10.1109/TIV.2024.3398129}{10.1109/TIV.2024.3398129}.

\bibitem{yangDataDrivenModelPredictive2025}
H.~Yang, Y.~He, C.~Mu, and P.~Li,
``Data-driven model predictive longitudinal control for autonomous
vehicles with unknown chassis,''
\emph{IEEE Trans. Veh. Technol.}, vol.~74, no.~8,
pp.~11732--11742, Aug. 2025, doi:
\href{https://doi.org/10.1109/TVT.2025.3552933}{10.1109/TVT.2025.3552933}.

\end{thebibliography}

% ======================================================================
\ifIncludeBiographies
% BIOGRAFIAS
% CONTENIDO: formación, puesto actual e intereses reales de cada autor.
% TONO/EXTENSION: un párrafo breve; unas 60--100 palabras orientativas.
% REVISION: coherencia con afiliaciones y autoría. T-IV pide biografías
% breves; incorporar fotografías reales cuando proceda en el paquete final.

\begin{IEEEbiographynophoto}{First A. Author}
\placeholder{Insert verified education, current position, and research
interests. Repeat for every author.}
\end{IEEEbiographynophoto}

\begin{IEEEbiographynophoto}{Second B. Author}
\placeholder{Insert this author's verified biography.}
\end{IEEEbiographynophoto}

\begin{IEEEbiographynophoto}{Third C. Author}
\placeholder{Insert this author's verified biography.}
\end{IEEEbiographynophoto}
\fi

\end{document}